\section{Experiments and discussion}
\label{sec:experiments}

\subsection{Single-stream comparison (Task 4.1)}
All four models are trained from random initialization using Adam with
learning rate $10^{-3}$, seed 42, and a maximum of 50 epochs. The cluster
runs use PyTorch 2.8.0 with CUDA on an NVIDIA A100 GPU partition
(MIG 1g.10gb). Early stopping monitors validation video loss with patience
five and a minimum improvement of $10^{-4}$. Stopping and checkpoint
selection are distinct: we report the highest-accuracy checkpoint, with
lower validation loss breaking ties, rather than automatically taking the
last epoch or the lowest-loss checkpoint.

Per-frame training uses batches of eight frames, giving 625 optimizer
updates per epoch from 5,000 frames. Late and early fusion use batches of
two videos (250 updates per epoch), while the 3D CNN uses one video per
batch (500 updates). Validation uses one video per batch for every model.
Per-frame, early fusion, and the 3D CNN stop after 31, 18, and 23 epochs;
late fusion reaches the 50-epoch limit. Total updates are therefore 19,375,
12,500, 4,500, and 11,500 for per-frame, late fusion, early fusion, and 3D
CNN, respectively. This is a comparison under a shared stopping rule,
not equal optimization or computation budgets.

% Generated by build_assets.py from explicit completed cluster runs.
\begin{table}[t]
    \centering
    \caption{Original-split validation results from the cluster runs.
    Epochs denotes selected/completed epochs. Accuracy and loss come
    from the same accuracy-selected checkpoint, not the final epoch.
    No test results are included.}
    \label{tab:baselines}
    \begin{tabular}{lrrr}
        \toprule
        Model & Epochs & Val. (\%) & Val. loss \\
        \midrule
        Per-frame & 22/31 & 90.8 & 0.567 \\
        Late fusion & 46/50 & 75.0 & 0.969 \\
        Early fusion & 18/18 & 22.5 & 2.104 \\
        3D CNN & 22/23 & 32.5 & 1.950 \\
        \bottomrule
    \end{tabular}
\end{table}


Per-frame aggregation achieves the highest validation accuracy at 90.8\%
(109/120 videos), followed by late fusion at 75.0\% (90/120), the 3D CNN
at 32.5\% (39/120), and early fusion at 22.5\% (27/120).
Table~\ref{tab:baselines} reports accuracy and loss at the same selected
checkpoint. For example, per-frame accuracy peaks at epoch 22, whereas its
lowest validation loss is 0.422 at epoch 26, with 88.3\% accuracy. Thus,
minimum loss and maximum accuracy need not identify the same epoch.

\paragraph{Learning curves.}
Figure~\ref{fig:learning} places all training and validation losses on one
set of axes. Per-frame loss falls fastest, but validation loss fluctuates
above its epoch-26 minimum before stopping. Late fusion improves over a
longer period: its validation loss falls from 2.295 to 0.969 at epoch 46,
then rises to 1.145 by epoch 50 while training loss continues to decrease.
This late divergence is consistent with some overfitting. Early fusion and
the 3D CNN finish with training losses of 1.907 and 1.633 and training
accuracies of 29.0\% and 41.6\%, respectively. Their weak validation
performance coincides with incomplete training-set fit; these runs do not
establish that the architectures lack useful temporal information. In
particular, loss-based early stopping can halt training despite continued
improvement in accuracy, as in early fusion.

Per-frame training loss is frame-level cross-entropy, whereas validation
loss is video-level negative log-likelihood after probability averaging.
Their gap is not a like-for-like generalization gap. The other models use
video cross-entropy for both splits. Training losses average batches during
weight updates; validation uses the end-of-epoch model. All results use one
seed and the original splits with subject overlap. They identify per-frame
aggregation as the current candidate for the corrected-split experiment,
but cannot establish its performance on unseen subjects. Test results
remain pending.

\subsection{Corrected-split comparison (Task 4.2.1)}
\label{sec:corrected}
\textbf{Draft status: corrected-split results are pending.}
Per-frame aggregation is selected using the original-split validation
results. Retrain it from random initialization
on the supplied corrected training split, keeping the architecture,
preprocessing, learning rate, seed, and training budget matched to the
original-split experiment. Use each regime's validation set for checkpoint
selection, then evaluate the selected checkpoints on their respective test
sets. Record the corrected split sizes and report accuracy changes in
percentage points. If the training recipe changes, rerun the original
baseline under that recipe too.

A drop would be consistent with optimistic evaluation under subject
overlap, but its magnitude must be measured. The split populations also
change, so the comparison is not a paired experiment on identical test
videos. No corrected-split score or leakage-related drop is claimed here.

\subsection{RGB and optical flow (Task 4.2.2)}
\label{sec:dualstream}
\textbf{Draft status: dual-stream implementation and results are pending.}
Two-stream networks combine an appearance stream with a motion stream
based on optical flow~\cite{simonyan2014twostream}. The corrected course
dataset supplies flow generated using RAFT~\cite{course2026,teed2020raft}.
Before finalizing the flow encoder, inspect the supplied representation,
its channel count, and its alignment with the RGB frames. In particular,
colored flow PNGs must not be described as raw horizontal and vertical
displacement fields without verifying their encoding.

The remaining experiment will compare RGB-only, flow-only, and fused
predictions on the corrected splits. Document the actual encoders,
normalization, fusion operation, and training schedule, and select any
fusion weight using validation data. Train any source encoder on corrected
training data as well. Report final test results only after model selection;
the current draft makes no claim that optical flow improves accuracy.
